%
% ungerade.tex -- Graph der Funktion f(x) = x/(1+x^2)
%
% (c) 2026 Damien Flury, OST Ostschweizer Fachhochschule
%
\documentclass[tikz]{standalone}
\usepackage{amsmath}
\usepackage{times}
\usepackage{txfonts}
\usepackage{pgfplots}
\usetikzlibrary{arrows,math}
\definecolor{darkgreen}{rgb}{0,0.6,0}
\begin{document}
\def\skala{1}
\begin{tikzpicture}[>=latex,thick,scale=\skala,yscale=3]

% Flächen unter dem Integranden, symmetrisch um den Ursprung
\fill[darkgreen!30,opacity=0.5]
	(-4,0) -- plot[domain=-4:0,samples=200] ({\x},{\x/(1+\x*\x)}) -- cycle;
\fill[darkgreen!30,opacity=0.5]
	(0,0) -- plot[domain=0:4,samples=200] ({\x},{\x/(1+\x*\x)}) -- (4,0) -- cycle;

% Achsen
\draw[->] (-4.3,0) -- (4.5,0) coordinate[label={$x$}];
\draw[->] (0,-0.8) -- (0,0.9) coordinate[label={right:$y$}];

% Achsenbeschriftung
\foreach \x in {-4,-3,-2,-1,1,2,3,4}{
	\draw (\x,-0.017) -- (\x,0.017);
}
\foreach \x in {-4,-3,-2,-1}{
	\node at (\x,0) [above] {$\x$};
}
\foreach \x in {1,2,3,4}{
	\node at (\x,0) [below] {$\x$};
}
\draw (-0.05,0.5) -- (0.05,0.5);
\node at (0,0.5) [above left] {$\frac{1}{2}$};
\draw (-0.05,-0.5) -- (0.05,-0.5);
\node at (0,-0.5) [below right] {$-\frac{1}{2}$};

% Graph
\draw[color=darkgreen,line width=1.4pt]
	plot[domain=-4.2:4.2,samples=200] ({\x},{\x/(1+\x*\x)});

\node[color=darkgreen] at (3.0,0.55) {$f(x)=\dfrac{x}{1+x^2}$};

\end{tikzpicture}
\end{document}
